\section{Finiteness of approximate representation matches}
\label{sec:nuclear}

We prove that, for a fixed unitary tuple $\mathbf u$ in a nuclear algebra,
only finitely many finite-dimensional irreducible representation classes
satisfy $d(\mathbf u,\pi)<a$.
The tuple need not satisfy any group relations. This distinction will allow
us to apply the result to coordinate lifts from a norm ultrapower.

We retain the notation and constants of Section~\ref{sec:spectral}.
For $C^*$-algebras $B$ and $C$, the minimal tensor product
$B\otimes_{\min}C$ is realized by faithful representations on separate
Hilbert spaces; the maximal tensor product $B\otimes_{\max}C$ is
characterized by pairs of commuting representations on one Hilbert space.
We use the tensor-norm characterization of nuclearity developed by
Takesaki and Lance~\cite{Takesaki1964,Lance1973}: $B$ is nuclear if
these two tensor norms agree for every $C$.
We first obtain a lower bound when the second tensor factor has no
finite-dimensional subrepresentations. The passage from a vector to its
reduced density operator removes the arbitrary unitaries in the first
factor.

\begin{lemma}[Square-root trace inequality]\label{lem:square-root-trace}
For positive Hilbert--Schmidt operators $X,Y$ on any Hilbert space,
\[
 \|X-Y\|_2^2\leq\|X^2-Y^2\|_1.
\]
Here $\|\cdot\|_1$ and $\|\cdot\|_2$ denote the trace and
Hilbert--Schmidt norms, respectively.
\end{lemma}

\begin{proof}
This is the square-root trace inequality of Powers--St\o rmer
\cite[Lemma~4.1]{PowersStormer1970};
we include the short argument needed here.
Put $F=X-Y$, and let $J=\operatorname{sgn}(F)$, with value zero on
$\ker F$. Since $X^2-Y^2=FX+YF$ and $JF=FJ=|F|$,
\begin{align*}
 \|X^2-Y^2\|_1
 &\geq \operatorname{Tr}\bigl(J(X^2-Y^2)\bigr)\\
 &=\operatorname{Tr}\bigl(|F|(X+Y)\bigr)\\
 &\geq\operatorname{Tr}(F_+F)+\operatorname{Tr}(F_-(-F))
 =\operatorname{Tr}(F^2).
\end{align*}
For the last inequality, use $X+Y-F=2Y\geq0$ and
$X+Y+F=2X\geq0$, and pair these inequalities with $F_+$ and $F_-$,
respectively.

All these traces are well defined in infinite dimension. Products of two
Hilbert--Schmidt operators are trace class; their traces are cyclic, by
finite-rank approximation and
$\|AB\|_1\leq\|A\|_2\|B\|_2$.
The trace of a product of two positive Hilbert--Schmidt operators is
nonnegative, as is seen by diagonalizing one factor. These observations
also justify the cyclic rearrangement with the bounded operator $J$.
The operators involved vanish on the orthogonal complement of a separable
subspace, so no separability hypothesis on the ambient Hilbert space is
needed.
\end{proof}

Bekka~\cite{Bekka1990} proved that if a unitary representation of a
property~(T) group on a Hilbert space $K$ admits a state on
$\mathcal B(K)$ invariant under conjugation by the representation,
then it has a nonzero finite-dimensional subrepresentation; see
also~\cite[proof of Proposition~2]{Valette1997}. The next lemma gives
the explicit spectral bound needed here for an arbitrary unitary tuple
in the first tensor factor.

\Needspace{7\baselineskip}
\begin{lemma}[Uniform spectral bound]
\label{lem:weak-mixing-bound}
Let $\Gamma,S,\kappa$ satisfy \eqref{eq:kazhdan-gap}. Let
$\sigma:\Gamma\to\mathcal U(K)$ be a unitary representation with no
nonzero finite-dimensional subrepresentation. For any Hilbert space $E$
and any unitaries $(U_s)_{s\in S}$ on $E$,
\[
 L\bigl((U_s\otimes\sigma(s))_{s\in S}\bigr)\geq b1,
 \qquad b=\frac{\kappa^2}{4|S|}.
\]
No relations among the $U_s$ are required.
\end{lemma}

\begin{proof}
Conjugation by $\sigma$ is a unitary representation on the Hilbert space
$\mathcal S_2(K)$ of Hilbert--Schmidt operators. It has no nonzero invariant
vector: if $A\neq0$ commutes with every $\sigma(g)$, then so does the
nonzero positive compact operator $|A|$. The eigenspace of any nonzero
eigenvalue of $|A|$ is a nonzero finite-dimensional $\sigma$-invariant
subspace.

Use inner products linear in the first variable. Fix a unit vector
$\xi\in E\otimes K$. Its reduced density operator
$T$ on $K$ is characterized by
\[
 \operatorname{Tr}(TC)=\langle(1\otimes C)\xi,\xi\rangle
 \qquad(C\in\mathcal B(K)).
\]
Concretely, write $\xi=\sum_i e_i\otimes v_i$ with $(e_i)$ orthonormal
and $\sum_i\|v_i\|^2=1$. Then
$T=\sum_i |v_i\rangle\langle v_i|$, with convergence in trace norm;
in particular, $T\geq0$ and $\operatorname{Tr}(T)=1$. Such an expansion
uses only countably many vectors, even if the Hilbert spaces are
nonseparable.

For $s\in S$, put $\xi_s=(U_s\otimes\sigma(s))\xi$ and
$\delta_s=\|\xi_s-\xi\|$. The reduced density operator of $\xi_s$ is
$T_s=\sigma(s)T\sigma(s)^*$. Testing the two vector states against
$1\otimes C$, for $\|C\|\leq1$, gives
\[
 \|T_s-T\|_1\leq2\delta_s.
\]
Lemma~\ref{lem:square-root-trace} therefore gives
\[
 \|\sigma(s)T^{1/2}\sigma(s)^*-T^{1/2}\|_2^2
 \leq\|T_s-T\|_1\leq2\delta_s.
\]
Since $\|T^{1/2}\|_2=1$, applying \eqref{eq:kazhdan-gap} to the
conjugation representation and then Cauchy--Schwarz yields
\[
 \kappa\leq2\sum_{s\in S}\delta_s
 \leq2\sqrt{|S|}\left(\sum_{s\in S}\delta_s^2\right)^{1/2}.
\]
Thus $\sum_s\delta_s^2\geq b$. Taking all unit vectors $\xi$ proves
the operator inequality; when $E\otimes K=\{0\}$ it is vacuous.
\end{proof}

The following quotient argument uses Wassermann's property~(T)
method~\cite{Wassermann1991}; see the exposition
in~\cite[Section~2.1]{ManuilovThomsen2007}. The preceding uniform bound
allows us to use an arbitrary unitary tuple in the first factor.
Nuclearity will transfer the bound to the coordinate quotient and
contradict the assumed small spectral minima.

\begin{lemma}[Nuclear finiteness]\label{lem:nuclear-finiteness}
Let $\Gamma,S,\kappa,R,b,a$ be as above. For every nonzero unital
nuclear $C^*$-algebra $B$ and every unitary $S$-tuple $\mathbf u$ in $B$,
the set
\[
 D(\mathbf u)=\{[\pi]:d(\mathbf u,\pi)<a\}
\]
is finite, where $\pi$ ranges over the nonzero finite-dimensional
irreducible unitary representations of $\Gamma$ and $[\pi]$ denotes its
unitary equivalence class. No group relations on $\mathbf u$ are assumed.
\end{lemma}

\begin{proof}
Suppose that $(\pi_j)_{j\geq1}$ are pairwise inequivalent representations
with $[\pi_j]\in D(\mathbf u)$. Write $M_j=\mathcal B(\bar V_{\pi_j})$, and form
\[
 Q=\frac{\prod_{j\geq1}M_j}{\bigoplus_{j\geq1}M_j},
 \qquad q(g)=[(\bar\pi_j(g))_j]\quad(g\in\Gamma).
\]
The product consists of bounded sequences, and the direct sum is the ideal
of sequences whose norms tend to zero in the ordinary sense. In
particular, this auxiliary quotient does not involve an ultrafilter.

Choose a faithful unital representation $\eta:Q\to\mathcal B(K)$, and
put $\sigma=\eta\circ q$. We first show that $\sigma$ has no nonzero
finite-dimensional subrepresentations. Fix a
finite-dimensional irreducible representation $\tau$ of $\Gamma$.
Under the identification
\[
 \bar V_{\pi_j}\otimes\bar V_\tau
 \simeq\mathcal S_2(V_\tau,\bar V_{\pi_j}),
\]
invariant vectors for $\bar\pi_j\otimes\bar\tau$ correspond to
intertwiners from $\tau$ to $\bar\pi_j$. By irreducibility and
inequivalence, such an intertwiner can exist for at most one $j$.
Consequently,
\[
 L\bigl((\bar\pi_j(s)\otimes\bar\tau(s))_{s\in S}\bigr)
 \geq\kappa1
\]
outside at most one coordinate. Fixed-size matrix amplification commutes
with this quotient, entry by entry:
\[
 Q\otimes\mathcal B(\bar V_\tau)
 \simeq
 \frac{\prod_j(M_j\otimes\mathcal B(\bar V_\tau))}
 {\bigoplus_j(M_j\otimes\mathcal B(\bar V_\tau))}.
\]
The exceptional coordinate disappears, and hence
\begin{equation}\label{eq:quotient-excludes-finite-parts}
 L\bigl((q(s)\otimes\bar\tau(s))_{s\in S}\bigr)\geq\kappa1.
\end{equation}
If $\sigma$ contained $\tau$, then
$\sigma\otimes\bar\tau$ would have an invariant vector, contradicting
\eqref{eq:quotient-excludes-finite-parts}. Every nonzero finite-dimensional
unitary representation contains an irreducible subrepresentation, so
$\sigma$ has none.

Represent $B$ faithfully and unitally on a Hilbert space $E$. By
Lemma~\ref{lem:weak-mixing-bound},
\begin{equation}\label{eq:spatial-lower-bound}
 L\bigl((u_s\otimes q(s))_{s\in S}\bigr)\geq b1
 \quad\text{in }B\otimes_{\min}Q.
\end{equation}
We now use nuclearity to pass this inequality to the coordinate quotient.
Let
\[
 C=\frac{\prod_j(B\otimes M_j)}{\bigoplus_j(B\otimes M_j)}.
\]
There are commuting unital $*$-homomorphisms
\[
 \alpha:B\longrightarrow C,\quad
 \alpha(x)=[(x\otimes1_{M_j})_j],
 \qquad
 \beta:Q\longrightarrow C,\quad
 \beta([(y_j)_j])=[(1_B\otimes y_j)_j].
\]
The map $\beta$ is well defined because
$\|1_B\otimes y_j\|=\|y_j\|$. The universal property of the maximal
tensor product gives a unital $*$-homomorphism
\[
 \Psi:B\otimes_{\max}Q\longrightarrow C,
 \qquad \Psi(x\otimes q)=\alpha(x)\beta(q).
\]
Since $B$ is nuclear,
$B\otimes_{\min}Q=B\otimes_{\max}Q$, so \eqref{eq:spatial-lower-bound}
passes through $\Psi$.

Put $L_j=L((u_s\otimes\bar\pi_j(s))_{s\in S})$. We have proved
$[(L_j)_j]\geq b1$ in $C$. Continuous functional calculus for
$f(t)=(b-t)_+$ now gives
\[
 [(f(L_j))_j]=f([(L_j)_j])=0,
 \qquad
 \|f(L_j)\|=(b-d(\mathbf u,\pi_j))_+\longrightarrow0.
\]
Thus $\liminf_j d(\mathbf u,\pi_j)\geq b$, whereas every term is
less than $a<b$. This contradiction proves finiteness.
\end{proof}

The only use of nuclearity was to identify the minimal and maximal tensor
products before applying $\Psi$. In particular, the proof does not assume
an identification of the minimal tensor product with the coordinate
quotient, and it imposes no bound on the dimensions of the $\pi_j$.
